% GENERATED FILE. Edit canonical lessons, then run npm run generate:books.
% canonical-source-hash: fnv1a64:a59ac6f40be489bb
% canonical-lessons: PA-C267-thana, PA-C267-manzil, PA-C267-dakkhana, PA-C267-adalat, PA-C267-thikana

\chapter{Which Floor, Which Office?}
\label{ch:pa-a2-which-floor-which-office}
\hlchaptermodality{267}

\begin{chapteropening}
\textbf{By the end of this chapter:} \emph{I can ask where an office is in a town.}

Complete the last lesson of chapter 267: I can ask where an office is in a town.
\end{chapteropening}

\section[\texorpdfstring{\pa{ਥਾਣਾ} (thāṇā)}{ਥਾਣਾ (thāṇā)}]{\pa{ਥਾਣਾ} (thāṇā) — a police station}

\label{lesson:PA-C267-thana}



\begin{quote}
Before the new one: say the Punjabi for a roadside eatery, then the Punjabi for a building.
\end{quote}

\subsection*{You'll want to know: \pa{ਥਾਣਾ}}
\textbf{\pa{ਥਾਣਾ}} — \emph{thāṇā} — \textquotedblleft{}a police station\textquotedblright{}.

\begin{quote}
\pa{ਥਾਣਾ} \pa{ਬਾਜ਼ਾਰ} \pa{ਦੇ} \pa{ਨੇੜੇ} \pa{ਹੈ।}
\end{quote}

\begin{quote}
\emph{thāṇā bāzār de neṛe hai.}
\end{quote}

\begin{quote}
The police station is near the market.
\end{quote}

A masculine noun.

\begin{cousinweb}
\textbf{\pa{ਥਾਣਾ}} (\emph{thāṇā}) is inherited from Sanskrit \emph{sthāna}, \textquotedblleft{}firm stance, place of abode\textquotedblright{}; CDIAL lists Punjabi \emph{thāṇā}, \textquotedblleft{}place, police station\textquotedblright{}, under it (CDIAL 13753).
\end{cousinweb}

\begin{grammarlens}[title={What you've built}]
One more word for this chapter.
\end{grammarlens}

\subsection*{Your turn}
\begin{itemize}
\raggedright
  \item \emph{Say it:} \emph{thāṇā}
  \item \emph{Say it:} \emph{thāṇā}, once more
  \item \emph{Say it:} \emph{thāṇā bāzār de neṛe hai.} — The police station is near the market.
  \item \emph{Recall:} say the Punjabi for a roadside eatery, then the Punjabi for a building, then say \emph{thāṇā} again
\end{itemize}

\begin{tcolorbox}[breakable,colback=teal!4,colframe=teal!35!black,title={Before you move on}]
What does \pa{ਥਾਣਾ} mean? (\textquotedblleft{}A police station\textquotedblright{}.) Say it once more.
\end{tcolorbox}
\section[\texorpdfstring{\pa{ਮੰਜ਼ਿਲ} (manzil)}{ਮੰਜ਼ਿਲ (manzil)}]{\pa{ਮੰਜ਼ਿਲ} (manzil) — a storey; a floor of a building}

\label{lesson:PA-C267-manzil}



\begin{quote}
Before the new one: say the Punjabi for a building, then the Punjabi for a police station.
\end{quote}

\subsection*{You'll want to know: \pa{ਮੰਜ਼ਿਲ}}
\textbf{\pa{ਮੰਜ਼ਿਲ}} — \emph{manzil} — \textquotedblleft{}a storey; a floor of a building\textquotedblright{}.

\begin{quote}
\pa{ਮੇਰਾ} \pa{ਦਫ਼ਤਰ} \pa{ਦੂਜੀ} \pa{ਮੰਜ਼ਿਲ} \pa{ਉੱਤੇ} \pa{ਹੈ।}
\end{quote}

\begin{quote}
\emph{merā daftar dūjī manzil utte hai.}
\end{quote}

\begin{quote}
My office is on the second floor.
\end{quote}

A feminine noun.

\begin{cousinweb}
\textbf{\pa{ਮੰਜ਼ਿਲ}} (\emph{manzil}) is a loan, usually traced through Persian to Arabic \emph{manzil}, \textquotedblleft{}a halting place\textquotedblright{}.
\end{cousinweb}

\begin{grammarlens}[title={What you've built}]
One more word for this chapter.
\end{grammarlens}

\subsection*{Your turn}
\begin{itemize}
\raggedright
  \item \emph{Say it:} \emph{manzil}
  \item \emph{Say it:} \emph{manzil}, once more
  \item \emph{Say it:} \emph{merā daftar dūjī manzil utte hai.} — My office is on the second floor.
  \item \emph{Recall:} say the Punjabi for a building, then the Punjabi for a police station, then say \emph{manzil} again
  \item \emph{Write it:} copy \textbf{\pa{ਮੰਜ਼ਿਲ}} once, checking the whole word
\end{itemize}

\begin{tcolorbox}[breakable,colback=teal!4,colframe=teal!35!black,title={Before you move on}]
What does \pa{ਮੰਜ਼ਿਲ} mean? (\textquotedblleft{}A storey; a floor of a building\textquotedblright{}.) Say it once more.
\end{tcolorbox}
\section[\texorpdfstring{\pa{ਡਾਕਖ਼ਾਨਾ} (ḍākkhānā)}{ਡਾਕਖ਼ਾਨਾ (ḍākkhānā)}]{\pa{ਡਾਕਖ਼ਾਨਾ} (ḍākkhānā) — a post office}

\label{lesson:PA-C267-dakkhana}



\begin{quote}
Before the new one: say the Punjabi for a police station, then the Punjabi for a storey.
\end{quote}

\subsection*{You'll want to know: \pa{ਡਾਕਖ਼ਾਨਾ}}
\textbf{\pa{ਡਾਕਖ਼ਾਨਾ}} — \emph{ḍākkhānā} — \textquotedblleft{}a post office\textquotedblright{}.

\begin{quote}
\pa{ਡਾਕਖ਼ਾਨਾ} \pa{ਥਾਣੇ} \pa{ਦੇ} \pa{ਸਾਹਮਣੇ} \pa{ਹੈ।}
\end{quote}

\begin{quote}
\emph{ḍākkhānā thāṇe de sāhmaṇe hai.}
\end{quote}

\begin{quote}
The post office is opposite the police station.
\end{quote}

A masculine noun.

\begin{cousinweb}
\textbf{\pa{ਡਾਕਖ਼ਾਨਾ}} (\emph{ḍākkhānā}) is built from \emph{ḍāk}, \textquotedblleft{}the post\textquotedblright{}, for which CDIAL gives no Punjabi form (CDIAL 5519), and \emph{khānā}, \textquotedblleft{}house\textquotedblright{}, usually traced to Persian \emph{khāna}.
\end{cousinweb}

\begin{grammarlens}[title={What you've built}]
One more word for this chapter.
\end{grammarlens}

\subsection*{Your turn}
\begin{itemize}
\raggedright
  \item \emph{Say it:} \emph{ḍākkhānā}
  \item \emph{Say it:} \emph{ḍākkhānā}, once more
  \item \emph{Say it:} \emph{ḍākkhānā thāṇe de sāhmaṇe hai.} — The post office is opposite the police station.
  \item \emph{Recall:} say the Punjabi for a police station, then the Punjabi for a storey, then say \emph{ḍākkhānā} again
  \item \emph{Write it:} copy \textbf{\pa{ਡਾਕਖ਼ਾਨਾ}} once, checking the whole word
\end{itemize}

\begin{tcolorbox}[breakable,colback=teal!4,colframe=teal!35!black,title={Before you move on}]
What does \pa{ਡਾਕਖ਼ਾਨਾ} mean? (\textquotedblleft{}A post office\textquotedblright{}.) Say it once more.
\end{tcolorbox}
\section[\texorpdfstring{\pa{ਅਦਾਲਤ} (adālat)}{ਅਦਾਲਤ (adālat)}]{\pa{ਅਦਾਲਤ} (adālat) — a court of law}

\label{lesson:PA-C267-adalat}



\begin{quote}
Before the new one: say the Punjabi for a storey, then the Punjabi for a post office.
\end{quote}

\subsection*{You'll want to know: \pa{ਅਦਾਲਤ}}
\textbf{\pa{ਅਦਾਲਤ}} — \emph{adālat} — \textquotedblleft{}a court of law\textquotedblright{}.

\begin{quote}
\pa{ਅਦਾਲਤ} \pa{ਸ਼ਹਿਰ} \pa{ਦੇ} \pa{ਵਿਚਕਾਰ} \pa{ਹੈ।}
\end{quote}

\begin{quote}
\emph{adālat shahir de vichkār hai.}
\end{quote}

\begin{quote}
The court is in the middle of the town.
\end{quote}

A feminine noun.

\begin{cousinweb}
\textbf{\pa{ਅਦਾਲਤ}} (\emph{adālat}) is a loan, usually traced through Persian to Arabic \emph{adāla}, \textquotedblleft{}justice\textquotedblright{}.
\end{cousinweb}

\begin{grammarlens}[title={What you've built}]
One more word for this chapter.
\end{grammarlens}

\subsection*{Your turn}
\begin{itemize}
\raggedright
  \item \emph{Say it:} \emph{adālat}
  \item \emph{Say it:} \emph{adālat}, once more
  \item \emph{Say it:} \emph{adālat shahir de vichkār hai.} — The court is in the middle of the town.
  \item \emph{Recall:} say the Punjabi for a storey, then the Punjabi for a post office, then say \emph{adālat} again
  \item \emph{Write it:} copy \textbf{\pa{ਅਦਾਲਤ}} once, checking the whole word
\end{itemize}

\begin{tcolorbox}[breakable,colback=teal!4,colframe=teal!35!black,title={Before you move on}]
What does \pa{ਅਦਾਲਤ} mean? (\textquotedblleft{}A court of law\textquotedblright{}.) Say it once more.
\end{tcolorbox}
\section[\texorpdfstring{\pa{ਠਿਕਾਣਾ} (ṭhikāṇā)}{ਠਿਕਾਣਾ (ṭhikāṇā)}]{\pa{ਠਿਕਾਣਾ} (ṭhikāṇā) — whereabouts; a place to stay}

\label{lesson:PA-C267-thikana}



\begin{quote}
Before the new one: say the Punjabi for a post office, then the Punjabi for a court of law.
\end{quote}

\subsection*{You'll want to know: \pa{ਠਿਕਾਣਾ}}
\textbf{\pa{ਠਿਕਾਣਾ}} — \emph{ṭhikāṇā} — \textquotedblleft{}whereabouts; a place to stay\textquotedblright{}.

\begin{quote}
\pa{ਸ਼ਹਿਰ} \pa{ਵਿੱਚ} \pa{ਮੇਰਾ} \pa{ਕੋਈ} \pa{ਠਿਕਾਣਾ} \pa{ਨਹੀਂ} \pa{ਹੈ।}
\end{quote}

\begin{quote}
\emph{shahir vich merā koī ṭhikāṇā nahīṁ hai.}
\end{quote}

\begin{quote}
I have no place to stay in the city.
\end{quote}

A masculine noun.

\begin{cousinweb}
\textbf{\pa{ਠਿਕਾਣਾ}} (\emph{ṭhikāṇā}) is inherited from a reconstructed \emph{*ṭhīkka}, \textquotedblleft{}firm, right\textquotedblright{}; CDIAL lists Punjabi \emph{ṭhikāṇā}, \textquotedblleft{}station\textquotedblright{}, under it (CDIAL 5503).
\end{cousinweb}

\begin{grammarlens}[title={What you've built}]
One more word for this chapter.
\end{grammarlens}

\subsection*{Your turn}
\begin{itemize}
\raggedright
  \item \emph{Say it:} \emph{ṭhikāṇā}
  \item \emph{Say it:} \emph{ṭhikāṇā}, once more
  \item \emph{Say it:} \emph{shahir vich merā koī ṭhikāṇā nahīṁ hai.} — I have no place to stay in the city.
  \item \emph{Recall:} say the Punjabi for a post office, then the Punjabi for a court of law, then say \emph{ṭhikāṇā} again
  \item \emph{Write it:} copy \textbf{\pa{ਠਿਕਾਣਾ}} once, checking the whole word
\end{itemize}

\begin{tcolorbox}[breakable,colback=teal!4,colframe=teal!35!black,title={Before you move on}]
What does \pa{ਠਿਕਾਣਾ} mean? (\textquotedblleft{}Whereabouts; a place to stay\textquotedblright{}.) Say it once more.
\end{tcolorbox}
